\section{Introduction}

Suppose that two coefficient lanes contribute a local oriented covariance
$\ip{W}{B}$.  Knowing the two norms alone gives only Cauchy--Schwarz.  A more
informative interface fixes a unit direction $u$, records the moments
$b=\ip{u}{B}$ and $w=\ip{u}{W}$, and retains the remaining transverse energies.
The resulting geometry is elementary but dimension-sensitive: two transverse
directions fill a disk, whereas one transverse direction reaches only its
boundary.

This distinction matters in the present twin-prime structural program.
TPC-244 proved that a common outer block multiplier contributes through
$|C_h|^2\ip{w_h}{b_h}$, so its aggregate sign cannot control the same-block
main covariance.  The next local question is therefore whether longitudinal
information forces, permits, or excludes cancellation inside
$\ip{w_h}{b_h}$.  The exact answer is a translated disk, circle, singleton, or
empty set according to transverse dimension and realizability.

Our maximum claim is structural L1.  The theorem classifies all abstract
vectors with fixed data; it does not assert that the actual arithmetic vectors
range over that feasible set.  In particular, no committed source currently
defines the required one-dimensional block direction or attaches the literal
V59 lanes to the common coefficient space.
